\section{Whole-group tensor bound and simultaneous truncation}
\label{sec:peps_group_truncation}
The results below formalize the whole-group tensor bound
\cite[Lemma~6.2]{OpenAI2026PolynomialPEPS}, September 24, 2026,
\texttt{05-frames.tex}, lines 125--179, together with the two
dimension-free matrix estimates used in its proof. All tensors have finitely
many coordinates. For a matrix $M$, write $\lVert M\rVert_1$ for the sum of its
singular values and $\lVert M\rVert_2$ for its Hilbert--Schmidt norm.

\begin{definition}[Nuclear norm]
    \label{def:peps_group_nuclear_norm}
    \lean{Matrix.nuclearNorm}
    \leanok
    For a finite rectangular complex matrix $M$, the nuclear norm $\lVert M\rVert_1$
    is the sum of the singular values of $M$.
\end{definition}

\begin{lemma}[H\"older bound for a product with a transpose]
    \label{lem:peps_group_nuclear_holder}
    \lean{Matrix.nuclearNorm_mul_transpose_le}
    \leanok
    \uses{def:peps_group_nuclear_norm}
    For finite complex matrices $A$ and $B$ with the same column index set,
    $\lVert AB^{\mathsf T}\rVert_1\le\lVert A\rVert_2\lVert B\rVert_2$.
\end{lemma}

\begin{proof}\leanok
    Let $w_i$ be an orthonormal basis of right singular vectors of
    $M=AB^{\mathsf T}$, and let $u_i$ be the normalized images $Mw_i$ when
    these are nonzero. The $i$th singular value equals
    $\langle u_i, A(B^{\mathsf T}w_i)\rangle
        =\langle A^{*}u_i, B^{\mathsf T}w_i\rangle$, so it is at most
    $\lVert A^{*}u_i\rVert\,\lVert B^{\mathsf T}w_i\rVert$. Summing over $i$ and applying
    the Cauchy--Schwarz inequality gives
    $\lVert M\rVert_1\le(\sum_i\lVert A^{*}u_i\rVert^2)^{1/2}(\sum_i\lVert B^{\mathsf T}w_i\rVert^2)^{1/2}$.
    Since the $u_i$ and $w_i$ are orthonormal families, the two sums are at
    most $\lVert A\rVert_2^2$ and $\lVert B\rVert_2^2$.
\end{proof}

\begin{lemma}[Truncation to leading singular directions]
    \label{lem:peps_group_orthonormal_truncation}
    \lean{Matrix.exists_orthonormal_truncation, Matrix.orthonormalProjector}
    \leanok
    \uses{def:peps_group_nuclear_norm}
    Let $M$ be a finite complex matrix and $k\ge0$ an integer. There is an
    orthonormal family of at most $k$ vectors in the column space of
    coordinates of $M$ whose orthogonal projector $Q$ satisfies
    $(k+1)\lVert M-MQ\rVert_2^2\le\lVert M\rVert_1^2$. If the number of columns is at most
    $k$, then $MQ=M$.
\end{lemma}

\begin{proof}\leanok
    Take $Q$ to be the projector onto the first $\min(k,n)$ right singular
    vectors, where $n$ is the number of columns. Then $\lVert M-MQ\rVert_2^2$ is
    the sum $\sum_{j>k}s_j^2$ of the squares of the remaining singular values
    $s_1\ge s_2\ge\cdots\ge0$. Since the sequence decreases,
    $(k+1)s_{k+1}\le\sum_j s_j$ and $\sum_{j>k}s_j^2\le s_{k+1}\sum_j s_j$.
    Multiplying gives $(k+1)\sum_{j>k}s_j^2\le(\sum_j s_j)^2=\lVert M\rVert_1^2$.
\end{proof}

\begin{definition}[Coordinate contraction of a tensor graph]
    \label{def:peps_group_contraction}
    \lean{TNLean.PEPS.WholeGroup.contraction, TNLean.PEPS.WholeGroup.vertexNormSq}
    \leanok
    Let $(V,E)$ be a finite graph whose edges $e$ have endpoints
    $t(e),h(e)\in V$ and finite label sets $D_e$. At each vertex $v$ let
    $A_v$ be a tensor with one leg for each edge incident to $v$ and one
    group of open legs with finite configuration set $P_v$. The coordinate
    contraction is
    \[
        Z(\sigma)=\sum_{\beta}\prod_{v\in V}A_v\bigl(\beta|_{v},\sigma_v\bigr),
    \]
    where $\beta$ runs over labellings of all edges and $\beta|_v$ is the
    restriction of $\beta$ to the edges incident to $v$. The norm of $A_v$
    is its Hilbert-space norm over all its legs.
\end{definition}

\begin{theorem}[Whole-group tensor bound, norm part]
    \label{thm:peps_group_contraction_norm}
    \lean{TNLean.PEPS.WholeGroup.contraction_norm_le_one,
        TNLean.PEPS.WholeGroup.partialNormSq_le_one}
    \leanok
    \uses{def:peps_group_contraction}
    Suppose that no edge joins a vertex to itself and every vertex tensor
    has norm at most one. Then the contraction $Z$ has norm at most one.
    The same bound holds for the contraction of every vertex subset, with
    the open legs of the subset and the edges leaving it kept open.
\end{theorem}

\begin{proof}\leanok
    Argue by induction on the vertex subset, adding one vertex at a time.
    Merging two disjoint vertex sets contracts all edges between them at
    once. After grouping their common indices, the merged tensor is the
    product of the two partial contractions, read as matrices from the
    remaining legs to the common edges, up to a transpose. Hence its norm is
    at most the product of the two norms, by $\lVert AB\rVert_2\le\lVert A\rVert_2\lVert B\rVert_2$.
    A single vertex has the norm of its tensor, since it has no
    self-contraction. The empty set has norm one. Graphs with cycles are
    allowed.
\end{proof}

\begin{theorem}[Whole-group tensor bound, nuclear part]
    \label{thm:peps_group_contraction_nuclear}
    \lean{TNLean.PEPS.WholeGroup.nuclearNorm_flattening_le_one,
        TNLean.PEPS.WholeGroup.flattening,
        TNLean.PEPS.WholeGroup.flattening_contraction}
    \leanok
    \uses{def:peps_group_contraction, def:peps_group_nuclear_norm}
    Under the hypotheses of Theorem~\ref{thm:peps_group_contraction_norm},
    for every set $T$ of vertices, the matrix reshaping of $Z$ with rows the
    open-leg groups of $T$ and columns those of the complement has nuclear
    norm at most one.
\end{theorem}

\begin{proof}\leanok
    \uses{thm:peps_group_contraction_norm, lem:peps_group_nuclear_holder}
    Contract $T$ and its complement internally, giving tensors $A$ and $B$
    of norm at most one by Theorem~\ref{thm:peps_group_contraction_norm}.
    Group all edges crossing the bipartition into one index. The reshaping
    of $Z$ is then $AB^{\mathsf T}$, so Lemma~\ref{lem:peps_group_nuclear_holder}
    gives $\lVert AB^{\mathsf T}\rVert_1\le\lVert A\rVert_2\lVert B\rVert_2\le1$. No dimension
    enters.
\end{proof}

\begin{lemma}[Simultaneous truncation of open-leg groups]
    \label{lem:peps_group_truncation_abstract}
    \lean{TNLean.PEPS.WholeGroup.exists_group_truncation,
        TNLean.PEPS.WholeGroup.norm_sub_factorMatrix_le,
        TNLean.PEPS.WholeGroup.factorMatrix,
        TNLean.PEPS.WholeGroup.productVector}
    \leanok
    \uses{def:peps_group_nuclear_norm}
    Let $Z$ be a tensor on $\bigotimes_{v\in V}\mathbb C^{P_v}$ of norm at
    most one whose reshaping across each single factor $v$ and its
    complement has nuclear norm at most one. Let $G\subseteq V$ contain
    every $v$ with $|P_v|>1$, and let $k\ge1$. Then there are orthonormal
    families of at most $k$ vectors in the respective factors whose
    tensor-product projection sends $Z$ to $Z_k$ with
    $\lVert Z-Z_k\rVert\le|G|/\sqrt{k}$. In those families $Z_k$ has at most
    $k^{|G|}$ product terms, with coefficients of absolute value at most one.
\end{lemma}

\begin{proof}\leanok
    \uses{lem:peps_group_orthonormal_truncation}
    For each factor $v$, Lemma~\ref{lem:peps_group_orthonormal_truncation}
    applied to the single-factor reshaping gives a projection $Q_v$ of rank
    at most $k$ with $\lVert Z-Q_v Z\rVert\le1/\sqrt{k+1}\le1/\sqrt{k}$, and with
    $Q_v Z=Z$ when $|P_v|\le1$. The projections act on disjoint tensor
    factors, so they commute and each has operator norm at most one.
    Telescoping their product over $G$ bounds $\lVert Z-Z_k\rVert$ by the sum of
    the separate errors. The tensor $Z_k$ has norm at most one, so each of
    its coefficients in the orthonormal product basis has absolute value at
    most one.
\end{proof}

\begin{theorem}[Whole-group tensor bound]
    \label{thm:peps_group_tensor}
    \lean{TNLean.PEPS.WholeGroup.exists_simultaneous_openLeg_truncation,
        TNLean.PEPS.WholeGroup.exists_simultaneous_group_truncation}
    \leanok
    \uses{def:peps_group_contraction}
    Let a finite graph carry at each vertex a tensor of norm at most one,
    contract its internal edges by the coordinate pairing, and suppose there
    are no edges from a vertex to itself. Group the open legs according to
    their original vertices. The resulting tensor $Z$ has norm at most one
    and has nuclear norm at most one across every bipartition into whole
    open-leg groups. If there are $g$ nonempty open-leg groups, then for
    every integer $k\ge1$ there are subspaces of dimension at most $k$ in the
    respective group spaces whose tensor-product projection sends $Z$ to a
    tensor $Z_k$ with $\lVert Z-Z_k\rVert\le g/\sqrt{k}$. In orthonormal bases of
    those subspaces, $Z_k$ has at most $k^g$ product terms, with coefficients
    of absolute value at most one.
\end{theorem}

\begin{proof}\leanok
    \uses{thm:peps_group_contraction_norm, thm:peps_group_contraction_nuclear,
        lem:peps_group_truncation_abstract}
    The norm and nuclear bounds are
    Theorems~\ref{thm:peps_group_contraction_norm}
    and~\ref{thm:peps_group_contraction_nuclear}; in particular every
    single-group reshaping has nuclear norm at most one. A vertex with no
    open legs has a one-point configuration set, so the set of $g$ vertices
    owning a nonempty open-leg group contains every factor with more than
    one configuration. Lemma~\ref{lem:peps_group_truncation_abstract} then
    gives the truncation.
\end{proof}

\begin{corollary}[Whole-group tensor bound for networks with edge matrices]
    \label{cor:peps_group_tensor_network}
    \lean{TNLean.PEPS.WholeGroup.network_norm_le_one,
        TNLean.PEPS.WholeGroup.nuclearNorm_flattening_network_le_one,
        TNLean.PEPS.WholeGroup.exists_simultaneous_network_truncation,
        TNLean.PEPS.WholeGroup.network_one_eq_contraction}
    \leanok
    \uses{thm:peps_group_tensor}
    Consider a tensor network in which each vertex carries one leg for each
    endpoint incidence and each edge carries a matrix. If every edge matrix
    is the identity, no edge joins a vertex to itself, and every vertex
    tensor has norm at most one, then the conclusions of
    Theorem~\ref{thm:peps_group_tensor} hold for the network contraction.
\end{corollary}

\begin{proof}\leanok
    \uses{thm:peps_group_tensor}
    With identity edge matrices the network contraction is the coordinate
    contraction. Without edges from a vertex to itself, the endpoint
    incidences at a vertex are exactly its incident edges, so the vertex
    norms agree, and the whole-group tensor bound applies.
\end{proof}
